\documentclass[11pt]{article}
\usepackage[margin=1in]{geometry}
\usepackage{amsmath,amssymb,booktabs,array,microtype}
\usepackage[colorlinks=true,linkcolor=blue,urlcolor=blue,citecolor=blue]{hyperref}
\usepackage{fancyhdr}
\pagestyle{fancy}\fancyhf{}\fancyhead[L]{Jones polynomial: proof audit}\fancyhead[R]{31 October 2026}\fancyfoot[C]{\thepage}
\setlength{\headheight}{14pt}
\setlength{\parindent}{0pt}\setlength{\parskip}{7pt}
\title{Jones unknot detection:\\a concrete obstruction to the Carmi--Cohen proof}
\author{Research and mathematical audit prepared with Codex}
\date{31 October 2026}
\begin{document}
\maketitle
\begin{abstract}
The June 21, 2026 preprint by Avishy Carmi and Eliahu Cohen claims to prove that the ordinary Jones polynomial detects the unknot. Its Section 5 geometric package implies a degree bound independent of the length of an ordinary anti-parallel twist. This consequence fails on an explicit twist-knot family based at the trefoil: the proposed bound is 22, whereas a permitted length-23 state has a nonzero Jones--Vassiliev coefficient $b_{51}=-1$. The calculation uses the paper's normalization and can be checked directly with the skein relation. It invalidates the proof as submitted; it gives no counterexample to the Jones unknot conjecture itself.
\end{abstract}

\section*{Scope and literature status}
The object under discussion is a classical knot in $S^3$, with normalized ordinary Jones polynomial $V_U=1$. The question is whether $V_K=1$ implies $K=U$.

The live arXiv record for Carmi--Cohen \cite{CC} lists only v1, submitted June 21, 2026, and gives no journal reference. Searches conducted for this audit did not locate a named independent mathematical validation or published refutation. This search result should not be read as a proof that no private or unindexed review exists.

Manolescu's ICM 2026 survey \cite{Manolescu} still describes the ordinary Jones question as open. Its online publication postdates the preprint, but its manuscript may have been written earlier: it is an expert status reference, not a direct assessment of Carmi--Cohen.

Several established results must be kept distinct:
\begin{itemize}
\item Tuzun--Sikora \cite{TS} verified Jones unknot detection for all knots with at most 24 crossings. A finite verification is not a proof for arbitrary crossing number.
\item Kronheimer--Mrowka \cite{KM} proved that reduced Khovanov homology detects the unknot. Passing to the graded Euler characteristic can cancel homological information, so this does not prove the Jones assertion.
\item Jones detection holds for alternating knots, by the classical span theorem; adequate and semi-adequate cases and the history of these partial results are discussed in \cite{AL}.
\item Bigelow \cite{Bigelow} describes representation-theoretic routes to the problem. These illustrate that the general question reaches beyond the verified special classes.
\end{itemize}
The conclusion warranted by the sources located here is that there is no verified resolution to report, and the particular June 21 proof fails the explicit test below.

\section*{An explicit family requiring no conjectural knot}
Use the three-crossing trefoil diagram in the opening illustration of Section 5 of \cite{CC} (printed p.~23). Replace its circled crossing by an odd anti-parallel twist of length $L$. Let the resulting knot be $K_L$, and write $V_L=V_{K_L}$. Choose chirality so that
\[
V_1=x^2+x^6-x^8.
\]
Changing the original seed crossing unknots the trefoil. In the normalization
\[
x^{-2}V_+-x^2V_-=(x-x^{-1})V_0,\qquad p=x-x^{-1},
\]
the polynomial $H$ of the oriented smoothing is therefore determined by
\[
pH=x^{-2}(x^2+x^6-x^8)-x^2
=1-x^2+x^4-x^6.
\]
Thus
\[
H=-x-x^5.
\]

For any larger odd $L$, changing one crossing of the positive twist cancels a pair by Reidemeister II, leaving $K_{L-2}$. Smoothing that crossing turns the other twist crossings into removable kinks and leaves the same smoothing $H$. Consequently,
\begin{equation}\label{rec}
V_L=x^4V_{L-2}+px^2H
=x^4V_{L-2}+x^2-x^4+x^6-x^8.
\end{equation}
This is an ordinary exact skein computation. In particular,
\[
V_3=x^2-x^4+2x^6-x^8+x^{10}-x^{12}.
\]
It agrees, after $q=x^2$, with the published polynomial for $5_2$ in Garoufalidis--Zagier \cite{GZ}, Appendix A.2, p.~79. Their p.~74 also confirms the starting trefoil polynomial.

Equation \eqref{rec} immediately shows that the unique highest monomial is
\[
-x^{2L+6}.
\]
Indeed, the initial highest exponent is 8, each step raises it by 4, and the inhomogeneous summand has degree only 8.

\section*{Converting this into the paper's exact coefficients}
The paper packages the Jones polynomial uniquely in the form
\[
V_L=A_L(p)+B_L(p)x,\qquad x^2=px+1,
\]
with $A_L,B_L\in\mathbb Z[p]$. Write
\[
A_L=\sum_q a_q(K_L)p^q,\qquad
B_L=\sum_q b_q(K_L)p^q.
\]

For positive $m$ the reduction of $x^m$ has $B$-component of degree $m-1$, with leading coefficient 1. This follows by induction: multiplying $A+Bx$ by $x$ sends
\[
(A,B)\longmapsto (B,A+pB).
\]
The $A$-component has smaller degree. All monomials of $V_L$ have positive powers of $x$, and every monomial other than $-x^{2L+6}$ has strictly smaller exponent. None can cancel that monomial's leading $p$-coefficient. Hence
\begin{equation}\label{actual}
\deg_p V_L=2L+5,\qquad b_{2L+5}(K_L)=-1.
\end{equation}
Here $\deg_p(A+Bx)$ means $\max(\deg A,\deg B)$, exactly the coefficient degree used in the preprint.

For $L=23$, equation \eqref{actual} gives
\[
\boxed{\deg_p V_{23}=51,\qquad b_{51}(K_{23})=-1.}
\]
No numerical approximation or search for an unknown knot is involved.

\section*{The incompatible claim in the preprint}
The relevant statements in \cite{CC} are:
\begin{enumerate}
\item Proposition 5.12 (printed pp.~38--48) asserts that its local construction supplies machine hypotheses (C1)--(C4), with overhead $\kappa_{\rm mach}=11$.
\item Theorem 5.6 (p.~30) derives a signed-machine degree bound $\widehat B=B_{\rm host}+11$, independent of $L$.
\item Lemma 5.8 (pp.~31--33), using (C4) to identify the doubly smoothed machine with the ordinary residual twist, derives
\[
f_r=0\quad\text{for every }r>\widehat B+2.
\]
\end{enumerate}
These statements precede and do not require the vanishing-anchor hypothesis of Theorem 5.10. In particular, the host need not have Jones polynomial 1. They therefore apply, as asserted, to the trefoil host just used.

That host has $n=3$ crossings and one component, so the paper's own host bound is $B_{\rm host}=3n=9$. Its package predicts
\[
\boxed{\deg_p V_L\le 9+11+2=22\quad\text{for every allowed odd }L.}
\]
Taking $L=23$ contradicts $b_{51}(K_{23})=-1$. This length also satisfies $L>\widehat B+2$, although Lemma 5.8 itself does not need that additional restriction.

The conditional algebra of Lemma 5.8 is not the source of the problem. Its paired skein identity
\[
V_{-+}+V_{+-}=x^{-4}V_{++}+x^4V_{--}-p^2V_{00}
\]
is correct. Multiplication by $x^{\pm4}$ raises the $p$-degree by at most four; dividing the displayed identity by $p^2$ then gives precisely the stated residual bound if all signed machine states really have the asserted uniform bound. Therefore the signed-state compression and the ordinary-twist identification cannot simultaneously hold as claimed. The geometric package in Proposition 5.12 cannot supply all the hypotheses used to reach that conclusion.

This audit does not assert that the exceptional-state move (G2), considered in isolation, has been uniquely identified as the offending move. The explicit calculation already disproves the combined geometric claims needed by the proof. A correction must say which claimed reduction or identification fails and replace it with a valid argument.

Changing the number 11 to another fixed constant would not repair the issue: the actual degrees $2L+5$ are unbounded.

A second check avoids the precise handedness and coefficient calculation. Choose the all-same-sign ordinary twist that preserves the reduced alternating trefoil diagram. The resulting reduced alternating knot has $L+2$ crossings, so the alternating-span theorem gives Jones span $L+2$ in $t$, or $2L+4$ in $x=t^{1/2}$; see \cite{Kauffman}. On the other hand, a $p$-degree bound of 22 in the basis $\{1,x\}$ confines the Laurent support to $[-22,23]$, because $p^q=(x-x^{-1})^q$. Its span would be at most 45. For $L=23$ the actual span is 50. This gives a separate contradiction to the same asserted bound.

\section*{What checks and what remains unsupported}
The rank-two ring construction, finite-type divisibility argument, coefficient skein recursions, and graph-support propagation were checked algebraically. No fatal error was found in these portions. This is a targeted mathematical audit, not a formal certification of every sentence.

The final component-count contradiction in Section 6 is also compatible with the familiar order-two Jones invariant once its hypothesized vanishing on the two-crossing face is granted. The missing step is the construction of that vanishing face: its initialization uses the false uniform compression package.

The reproducible companion \texttt{verify\_twist\_obstruction.py} uses only Python's standard library and integer arithmetic. It verifies the recurrence for every odd $1\le L\le23$, checks the trefoil and $5_2$ initial values, computes the rank-two reduction, and independently reconstructs each Laurent polynomial by substituting $p=x-x^{-1}$. The output includes:
\[
\begin{array}{c|r|r|r}
L & \max\deg_x V_L & \deg_p V_L & b_{\deg_p V_L}\\ \hline
1&8&7&-1\\
3&12&11&-1\\
5&16&15&-1\\
23&52&51&-1
\end{array}
\]

This is an obstruction to the submitted proof, not a counterexample to unknot detection. Every knot used in the calculation has a nontrivial Jones polynomial. The conjecture may still be true; Carmi--Cohen v1 does not establish it.

\begin{thebibliography}{9}\small
\bibitem{CC} A. Carmi and E. Cohen, \emph{On trivial Jones--Vassiliev polynomials}, arXiv:2606.22410v1, submitted June 21, 2026. \url{https://arxiv.org/abs/2606.22410v1}
\bibitem{Manolescu} C. Manolescu, \emph{From Knots to Four-Manifolds}, ICM 2026, Vol.~2, pp.~154--179, especially Section 2. \url{https://doi.org/10.1137/25M1806703}
\bibitem{TS} R. E. Tuzun and A. S. Sikora, \emph{Verification of the Jones Unknot Conjecture Up to 24 Crossings}, J. Knot Theory Ramifications 30(3), 2150020 (2021). \url{https://arxiv.org/abs/2003.06724}
\bibitem{KM} P. B. Kronheimer and T. S. Mrowka, \emph{Khovanov homology is an unknot-detector}, Publ. Math. IHES 113 (2011), 97--208. \url{https://arxiv.org/abs/1005.4346}
\bibitem{AL} A. M. Lowrance and D. Spyropoulos, \emph{The Jones polynomial of an almost alternating link}, New York J. Math. 23 (2017), 1611--1639. \url{https://nyjm.albany.edu/j/2017/23-71.html}
\bibitem{Bigelow} S. Bigelow, \emph{Does the Jones polynomial detect the unknot?}, J. Knot Theory Ramifications 11 (2002), 493--505. \url{https://arxiv.org/abs/math/0012086}
\bibitem{GZ} S. Garoufalidis and D. Zagier, \emph{Knots, Perturbative Series and Quantum Modularity}, SIGMA 20 (2024), 055, especially pp.~74 and 79. \url{https://doi.org/10.3842/SIGMA.2024.055}
\bibitem{Kauffman} L. H. Kauffman, \emph{New Invariants in the Theory of Knots}, American Mathematical Monthly 95(3) (1988), 195--242, especially p.~212. \url{https://homepages.math.uic.edu/~kauffman/Bracket.pdf}
\end{thebibliography}
\end{document}
